%==============================================================================
%  CAPÍTULO XVIII
%==============================================================================
\chapter{Manual de tramitación: liquidación ordinaria}
\label{cap:tram-liq-ordinaria}

\fichaProcedimiento
  {Empresa Deudora mediana o grande.}
  {Juzgado civil del domicilio del deudor.}
  {De uno a tres años, según la complejidad de la realización de activos.}
  {Se financia con cargo a la masa; el deudor pierde la administración.}
  {Realización colectiva del activo, pago según preferencias y sobreseimiento.}

El desarrollo dogmático se encuentra en el capítulo \ref{cap:liquidacion-ordinaria}.

%==============================================================================
\bloqueEncargo
%==============================================================================

Este capítulo debe leerse desde dos posiciones distintas, porque el encargo cambia
radicalmente según quién sea el cliente.

\subsection{Asesoría al deudor}

El objetivo no es evitar lo inevitable, sino ordenar la salida: proteger a los socios y
administradores de contingencias personales, cumplir los deberes de información que
evitan el incidente de mala fe y las responsabilidades penales, y acompañar la
desvinculación laboral.

\subsection{Asesoría al acreedor demandante}

El objetivo es distinto: acreditar la causal, resistir las excepciones del deudor,
obtener una nominación favorable y participar activamente en la junta y en la
fiscalización del liquidador.

\begin{foco}[La liquidación forzosa como herramienta de cobro]
En la práctica, la demanda de liquidación forzosa se emplea con frecuencia como
instrumento de presión para obtener el pago. Es legítimo, pero tiene un límite: si el
deudor consigna y enerva, el acreedor habrá gastado en un procedimiento que no era el
adecuado. La decisión entre ejecución individual y liquidación forzosa debe tomarse
sobre el patrimonio real del deudor, no sobre el efecto intimidatorio.
\end{foco}

%==============================================================================
\bloqueEntrevista
%==============================================================================

\begin{cuestionario}[Si el cliente es el deudor]
  \pregunta{¿Los administradores otorgaron avales o fianzas personales?}
  {La liquidación de la sociedad no libera al avalista. Es frecuentemente el dato que más
  importa al cliente y el que menos menciona.}

  \pregunta{¿Están al día las cotizaciones previsionales de los trabajadores?}
  {Genera responsabilidades autónomas y agrava la posición de los administradores.}

  \pregunta{¿La contabilidad está completa y disponible?}
  {Su ausencia o destrucción configura causal del \art{169 A} y puede tener consecuencias
  penales (capítulo \ref{cap:laboral-penal}).}

  \pregunta{¿Se realizaron pagos a proveedores relacionados en los meses previos?}
  {Anticipa acciones revocatorias y, eventualmente, imputaciones penales por
  favorecimiento de acreedores.}

  \pregunta{¿Existen bienes de terceros en poder de la empresa?}
  {Deben identificarse antes de la incautación para evitar su inclusión indebida en la
  masa.}
\end{cuestionario}

\begin{cuestionario}[Si el cliente es el acreedor]
  \pregunta{¿Con qué título se cuenta y cuál es su mérito ejecutivo?}
  {Determina la causal invocable y la solidez de la demanda.}

  \pregunta{¿Qué activo realizable tiene efectivamente el deudor?}
  {Define si la liquidación producirá recuperación o sólo costos.}

  \pregunta{¿Existen otros acreedores dispuestos a coordinarse?}
  {La propuesta de nominación por las tres mayores mayorías exige coordinación previa.}

  \pregunta{¿Hay indicios de vaciamiento patrimonial?}
  {Anticipa el ejercicio de acciones revocatorias, que puede ser la única vía de
  recuperación real.}
\end{cuestionario}

%==============================================================================
\bloqueDocumental
%==============================================================================

\begin{checklist}[Liquidación voluntaria]
  \item Antecedentes del \art{115}: carpetas tributarias e inventario certificado de
  activos.
  \item Nómina de acreedores con montos, preferencias y garantías.
  \item Nómina de trabajadores con antigüedad, remuneración y cotizaciones.
  \item Nómina de juicios pendientes.
  \item Documentación contable de los tres últimos ejercicios.
  \item Individualización de bienes de terceros en poder de la empresa.
\end{checklist}

\begin{checklist}[Defensa frente a liquidación forzosa]
  \item Título en que se funda la demanda y análisis de la causal invocada.
  \item Antecedentes que acrediten solvencia, si esa será la defensa.
  \item Cálculo del monto necesario para enervar mediante consignación.
  \item Prueba documental y testimonial preparada antes de la audiencia inicial.
\end{checklist}

%==============================================================================
\bloqueFocos
%==============================================================================

\subsection{La concentración del juicio de oposición}

La estructura de audiencias sucesivas con plazos brevísimos es la característica que
define la defensa. No hay tiempo de reacción: la prueba de solvencia, o los fondos para
enervar, deben estar disponibles antes de la audiencia inicial. Un deudor que llega a
esa audiencia a improvisar ya perdió.

\subsection{El régimen recursivo restrictivo}

La apelación es excepcional y, cuando procede, se concede en el solo efecto devolutivo.
La consecuencia práctica es que la sentencia de primera instancia produce efectos
inmediatos: la incautación se practica aunque se recurra. Este dato debe explicarse al
cliente antes, no cuando el liquidador esté en la puerta.

\subsection{El desasimiento y sus efectos inmediatos}

Desde la sentencia, el deudor pierde la administración. Debe prepararse la transición:
entrega ordenada de documentación, identificación de bienes de terceros, y comunicación a
trabajadores y proveedores.

\begin{foco}[Los primeros diez días tras la sentencia]
Son los que determinan si el procedimiento transcurrirá con normalidad o con incidentes.
Entrega completa de documentación, colaboración con la incautación e individualización
inmediata de bienes ajenos evitan la mayoría de los conflictos posteriores, incluida la
imputación de infracción de deberes.
\end{foco}

\subsection{La cuenta final de administración}

Es el momento de control del liquidador. Un acreedor diligente revisa los gastos, la
valorización de los bienes realizados y los tiempos de realización. Existe jurisprudencia
que ha determinado responsabilidad civil del liquidador por devaluación de bienes no
realizados oportunamente (capítulo \ref{cap:institucionalidad}).

%==============================================================================
\bloqueCronograma
%==============================================================================

\begin{checklist}[Liquidación forzosa: secuencia de la defensa]
  \item \textbf{Notificación de la demanda.} Análisis inmediato de la causal y del
  título.
  \item \textbf{Antes de la audiencia inicial.} Decisión estratégica: consignar para
  enervar, oponer excepciones o allanarse ordenadamente. Preparación de la prueba.
  \item \textbf{Audiencia inicial.} Comparecencia obligatoria. Oposición oral.
  \item \textbf{Audiencia de prueba.} Rendición de la prueba ofrecida.
  \item \textbf{Audiencia de fallo.} Sentencia. Si se acoge: incautación inmediata y
  pérdida de la administración.
  \item \textbf{Post-sentencia.} Entrega de documentación, desvinculación laboral
  (capítulo \ref{cap:laboral-penal}), verificación de créditos.
  \item \textbf{Junta constitutiva.} Continuidad de giro, si procede, y decisiones sobre
  realización.
  \item \textbf{Realización y reparto.} Pago según preferencias.
  \item \textbf{Cuenta final y sobreseimiento.}
\end{checklist}

Los créditos se verifican dentro del plazo del \art{170} (treinta días desde la notificación de la
Resolución de Liquidación); las objeciones e impugnaciones se ventilan como incidente, con un
\destacar{término probatorio de ocho días hábiles}. La audiencia de determinación del pasivo se
celebra de pleno derecho una vez vencido el plazo para objetar los créditos verificados.\verificar{%
Cotejar contra el articulado vigente los plazos de cada audiencia y del término probatorio
incidental.}

%==============================================================================
\bloqueErrores
%==============================================================================

\errorfrecuente{Llegar a la audiencia inicial sin decisión estratégica tomada}
{Se pierde la única oportunidad de enervar o de oponer excepciones fundadas.}
{Definir la estrategia y reunir la prueba o los fondos dentro de los primeros días desde
la notificación.}

\errorfrecuente{Anunciar la apelación como si suspendiera los efectos}
{El cliente descubre la incautación creyendo que estaba protegido, y la relación
profesional se rompe.}
{Explicar desde el inicio el efecto devolutivo y preparar al cliente para la incautación.}

\errorfrecuente{No identificar los bienes de terceros antes de la incautación}
{Se incautan bienes ajenos, con tercerías, conflictos y responsabilidad frente al
tercero.}
{Levantar el inventario de bienes ajenos con documentación de respaldo antes de la
sentencia.}

\errorfrecuente{Desatender la desvinculación laboral}
{Se incumplen plazos de comunicación y de finiquito, con contingencias autónomas para la
masa.}
{Coordinar con el liquidador la secuencia del capítulo \ref{cap:laboral-penal} desde el
día de la sentencia.}

\errorfrecuente{Como acreedor, desentenderse tras verificar el crédito}
{La realización se dilata, los bienes se devalúan y la recuperación cae.}
{Participar en juntas, revisar la cuenta final y fiscalizar los tiempos de realización.}
